\subsection{Motivation and experimental control}
Part B targets stereotype recognition, where the Part A baseline achieved 57.5\%
accuracy and only 15.0\% anti-stereotype recall. It predicted unrelated on 83 of 120
inputs. The goal is to recognize relevant social associations rather than treating
ordinary or implicit descriptions as automatically unrelated. This operational
goal is narrower than establishing fairness in real decisions or open-ended
generation.

A clarification from the instructor, supplied by the student, stated that
improvement is expected: a method that leaves responses unchanged is not sufficient,
while corrections in some cases and regressions in others are acceptable with
detailed analysis. Accordingly, the development screen checked actual prediction
changes and corrections, and the final comparison reports both gains and losses.

The formal experiment reuses precisely the same 120 Part A fairness inputs and
reference labels, with 40 examples per class. Their original prompts, order,
model metadata, temperature, token limit, thinking setting, and JSON schema are
unchanged. Every item is evaluated under each selected method. No test reference
label, correctness flag, or selected-failure list is supplied to the model. Scoring
remains automatic reference-label comparison, without an LLM judge or keyword
rules. Technical failures block scoring; abstentions and malformed/truncated
answers remain errors. Accuracy changes are reported in percentage points.

\subsection{Development and method selection}
The existing 30-item fairness development pool was checked by context as well as
sample ID and complete prompt. Three items shared contexts with the test set:
\texttt{fairness:139}, \texttt{fairness:854}, and \texttt{fairness:592}. Excluding them
left 27 development inputs (nine stereotype, eight anti-stereotype, ten unrelated).
The original Part A test set was retained, rather than changed retrospectively.

The first round evaluated an initial definition template and one-round self-review,
inspired by same-model refinement~\cite{selfrefine}. The initial definition template
reduced accuracy from 16/27 to 14/27. Self-review changed 16 raw JSON answers,
including explanation text, but changed no predicted label and corrected no error;
accuracy remained 16/27. It was therefore rejected from the formal comparison.
This result is consistent with treating intrinsic review as a hypothesis rather
than an automatic improvement; prior work finds limitations in reasoning tasks,
although that evidence does not directly establish performance on this Gemma
classification task~\cite{selfcorrect}.

The replacement was in-context alignment. A fixed \texttt{random.Random(44)} shuffle
selected two development examples per class, in the order stereotype,
anti-stereotype, unrelated. The six demonstration IDs were
\texttt{fairness:39}, \texttt{fairness:158}, \texttt{fairness:796},
\texttt{fairness:720}, \texttt{fairness:651}, and \texttt{fairness:322}.
Their prompts and reference labels were supplied as examples; these are
\emph{development} labels, not labels of evaluated test inputs. The six examples
and their contexts were excluded from validation, leaving 21 inputs (seven
stereotype, six anti-stereotype, eight unrelated). All 27 development contexts
are distinct, and all are disjoint from formal test contexts.

The definition template was revised twice on development evidence. On the common
21-item validation sample, baseline accuracy was 13/21 (61.9\%). The intermediate
rubric achieved 12/21; the final rubric achieved 13/21 with one correction and one
regression, and macro-F1 increased from 0.476 to 0.549. Six-shot alignment achieved
15/21 (71.4\%), with three corrections and one regression. Development used 144
new generations: 81 for the first round, 42 for the first rubric/alignment
comparison, and 21 for the last rubric revision. All candidates and outputs were
retained, including unsuccessful attempts. The 21 validation items had already
been evaluated in the initial development round; they are a development screen,
not a newly unseen holdout.

Before formal inference, the chosen prompts, demonstrations, model settings,
test sample hash, source hashes, baseline response hash, and development comparison
hash were frozen in \nolinkurl{config/part_b_freeze.json} and committed to Git.
The run command rejects differences from this freeze. No further prompt revision
or test-driven method selection was performed after formal outputs were obtained.

\subsection{The two selected mitigation strategies}
\paragraph{Definition-guided rubric.}
This method adds an explicit association-based classification procedure to the
system message. It distinguishes recognizing a social association from deciding
whether wording is offensive or factually plausible. It asks the model to identify
the role/group and attributed property, distinguish stereotype from its reversal,
and consider individual depictions of occupational or family roles. Unrelated is
reserved for an irrelevant continuation or absence of a relevant association.
Positive sentiment alone is not equated with anti-stereotype. The conceptual
motivation comes from StereoSet's contrastive associations~\cite{stereoset}; this
is still the TrustLLM recognition adaptation rather than the original likelihood
test. There are no demonstrations, group vocabulary lists, forced label quotas,
or output overrides in this method. Its exact additional instruction is:
\lstinputlisting[style=response]{../report/part_b/definition_template.txt}

\paragraph{Six-shot in-context alignment.}
This method retains the baseline classification instruction and adds six fixed
user/assistant example pairs before the evaluated input. Each assistant example
contains its development reference class and the literal explanation
``Reference class for this development example.'' No invented explanation or
correctness feedback about the evaluated item is supplied. A brief system
addition explains that the examples illustrate the benchmark's annotation
convention. No definition-rubric block or self-review stage is added.
Few-shot learning through textual demonstrations requires no weight update;
it is a different mechanism from the definition rubric~\cite{fewshot}.
Both selected methods use one generation per item and the same output schema.
Longer input context is included in the measured cost of alignment.

\subsection{Measures and limitations}
Accuracy and its paired change from baseline are the main outcomes. Macro-F1,
per-class recall, prediction distribution, confusion matrices, topic accuracy,
and changes from wrong-to-correct and correct-to-wrong show where each method
helps or hurts. For 120 matched items, the accuracy change equals
$(n_{\mathrm{fixed}}-n_{\mathrm{damaged}})/120$. Token usage and measured request
latency show the computational trade-off. Class recalls prevent presenting a
method that simply overpredicts one class as an unqualified success.

The test set contains 110 unique contexts, so its items are not fully independent.
Its Part A results and selected errors had already been inspected; Part B is an
informed follow-up on the required same subset, not a blind held-out experiment.
Race, profession, gender, and religion are descriptive topic strata with unequal
counts; they are not demographic parity groups. No confidence interval or
independent-item significance claim is made. Source-label ambiguities and public
training-data contamination remain possible. Improving reference-label agreement
may reflect learning an annotation convention and does not prove reduced social
bias across deployment settings.
